\section{Cấu hình huấn luyện}
\begin{table}[H]\centering\small
\caption{Cấu hình chính của ba thí nghiệm}\label{tab:traincfg}
\begin{tabular}{p{3.1cm}p{3.3cm}p{3.3cm}p{4.5cm}}\toprule
Thuộc tính & E3 YOLO11n & E4 YOLO11s & E6 Faster R-CNN\\\midrule
Khởi tạo & COCO pretrained & COCO pretrained & COCO V2 pretrained\\
Epoch & 10/30 lịch LR & 10/30 lịch LR & 10 theo code/checkpoint\\
Batch & 8 & 8 & 2\\
Kích thước & 640×640 & 640×640 & min 640, max 960\\
Optimizer & SGD & SGD & SGD\\
LR ban đầu & 0,01 & 0,01 & 0,005\\
Momentum & 0,937 & 0,937 & 0,9\\
Weight decay & 0,0005 & 0,0005 & 0,0005\\
Scheduler & Ultralytics 30 epoch & Ultralytics 30 epoch & StepLR(3; 0,1)\\
Seed & 0 & 0 & 0\\\bottomrule
\end{tabular}
\end{table}

E3 và E4 có 10 hàng lịch sử liên tục. Ở epoch 10, learning rate là 0,00703, còn cao hơn đích $lr_0\times lrf=0,0001$ của lịch 30 epoch. \texttt{close\_mosaic=10} chỉ đóng mosaic ở 10 epoch cuối; run dừng trước epoch 20 nên pha này chưa xảy ra. Hai run là 10 epoch đầu của lịch 30, không phải lịch 10 epoch hoàn chỉnh và không chứng minh hội tụ.

E6 dùng AMP, lưu trạng thái giữa epoch và backup định kỳ. Mỗi epoch đánh giá toàn VALID bằng TorchMetrics, bước scheduler rồi cập nhật best. Tuy nhiên \texttt{metrics.csv}, \texttt{config.json}, \texttt{environment.json} và \texttt{last\_hf\_manifest.json} rời thiếu ở local. Code và checkpoint chứng minh kiến trúc nạp được để suy luận; chúng không đủ xác minh độc lập GPU, phiên bản framework và đường cong từng epoch của run.

\section{Tính toàn vẹn checkpoint}
Ba checkpoint có SHA256 đã khóa: E3 \texttt{41f794...ed02}, E4 \texttt{dc09c7...e2f5}, E6 \texttt{ce459c...f9bd}. Adapter từ chối chạy khi hash sai. E6 mang đuôi \texttt{.zip} nhưng là state dictionary PyTorch, không phải archive; model được dựng lại bằng builder V2 với head 26 lớp và nạp strict.

Best checkpoint E3 có 5.477.402 byte, E4 19.191.770 byte, E6 173.898.073 byte. Kích thước tệp không bằng bộ nhớ runtime vì activation và workspace phụ thuộc ảnh, framework và thiết bị. Nó vẫn phản ánh chi phí lưu trữ và góp phần vào thời gian request đầu khi model chưa cache.

\section{Giao thức đánh giá chung}
Ground truth YOLO được chuyển từ tâm--rộng--cao chuẩn hóa sang hai góc pixel. Prediction E6 đổi foreground 1--25 về ROD 0--24; background bị bỏ. AP dùng 101 điểm recall tại IoU 0,50 đến 0,95, bước 0,05 và tối đa 100 detection/ảnh. Precision/Recall vận hành ghép tại IoU 0,50 với tối đa 300 detection/ảnh.

ROD không có nhãn crowd/ignore nên evaluator đặt area=all và không bỏ qua crowd. “COCO-style” chỉ mô tả dải IoU và nội suy; implementation không phải pycocotools chính thức. Ba JSON kết quả lưu protocol bên cạnh metric để số đo không bị tách khỏi định nghĩa.

\section{Khóa TEST và nguyên tắc lựa chọn}
Sau VALID, \texttt{evaluation\_lock.json} ghi hash data YAML, confidence config, checkpoint và protocol. Script TEST yêu cầu khóa hợp lệ cùng cờ xác nhận. Nếu model, threshold hoặc protocol đổi sau khi xem TEST, kết quả không còn độc lập và phải được coi là thí nghiệm khác.

TEST chưa chạy. Giữ trống có chủ đích tốt hơn điền metric từ VALID hoặc log framework. Khi đánh giá cuối, cả ba model cần chạy đúng một lượt trên 1.611 ảnh TEST, không chọn lại ngưỡng và báo metric tổng lẫn theo lớp.

\section{Thiết kế ứng dụng}
Gradio nhận JPG/JPEG/PNG, cho chọn E3/E4/E6, chạy adapter, vẽ bbox và trả PNG cùng danh sách detection. Adapter chuẩn hóa class ID, tọa độ, score và thời gian. Model được cache tối đa ba instance; CUDA được chọn khi có, nếu không dùng CPU. Đường dẫn checkpoint cố định và hash được kiểm trước khi load.

Thời gian suy luận đo riêng quanh model. Thời gian callback bắt đầu từ đối tượng PIL và kết thúc sau kiểm ảnh, inference, vẽ, lưu PNG và tạo văn bản. Nó không gồm upload qua mạng, hàng đợi server và browser render. Tách hai phạm vi đo giúp tránh gọi inference latency là trải nghiệm end-to-end.
